\documentclass[11pt,a4paper]{article}

% =========================
% Layout and fonts
% =========================
\usepackage[a4paper, margin=2cm]{geometry}

\usepackage{xetexko}
\usepackage{fontspec}
\setmainfont{TeX Gyre Pagella}
\setmainhangulfont{Noto Sans CJK KR}

\usepackage{setspace}
\singlespacing

\usepackage{indentfirst}
\setlength{\parindent}{10pt}

% =========================
% Tables and colors
% =========================
\usepackage{booktabs}
\usepackage{array}
\usepackage{colortbl}
\usepackage[dvipsnames]{xcolor}

% =========================
% Links
% =========================
\usepackage[
    colorlinks=true,
    linkcolor=MidnightBlue
]{hyperref}

% =========================
% Title
% =========================
\title{\bfseries 대내외 환경변화에 대응한 公社 역할 정립과 발전 방안}
\author{녹색전환 분야 연구 수행 과정 기록}
\date{\today}

% =========================
\begin{document}
% =========================

\maketitle

\section*{연구 개요}

본 연구의 녹색전환 파트는 기후위기 대응, 탄소중립, 자원순환 및 생태보전이라는
세 가지 축을 중심으로 농업ㆍ농촌 부문의 구조적 변화를 분석하고,
이를 한국농어촌공사의 기능 재정립 및 실천과제 도출로 연결하는 것을 목적으로 한다.

최근 기후위기 심화와 국제 식량시장 불안정성 증가는 농업ㆍ농촌의 생산 기반과
공공 인프라 운영 체계 전반에 구조적 변화를 요구하고 있다.
특히 농업용수 관리, 농지 정비, 재해 대응, 농촌공간 관리 기능을 수행하는
한국농어촌공사는 기존의 생산기반 유지ㆍ관리 기능을 넘어,
기후적응형 자원관리, 탄소중립형 기반시설 운영,
친환경 농촌공간 조성 기능까지 포괄하는 방향으로
역할 확대가 필요한 상황이다.

분석의 기본 흐름은 다음과 같다.

\begin{center}
\textbf{정책 변화 $\rightarrow$ 농업ㆍ농촌 영향 $\rightarrow$ 공사 역할 변화}
\end{center}

\section*{1단계: 연구 프레임 설정}

\subsection*{목적}

1단계에서는 녹색전환의 분석 범위를 명확히 설정하고,
이후 문헌조사, 정책분석, 사례조사, 공사 기능 재정립을 위한 기본 분석틀을 마련한다.

녹색전환은 단순히 온실가스 감축 정책에 한정되지 않으며,
기후위기에 대한 적응(adaptation),
온실가스 감축 및 탄소중립(mitigation),
자원순환 및 생태보전(resource circulation and ecological conservation)을
포괄하는 개념으로 정의한다.

\subsection*{주요 분석 범위}

\begin{table}[h]
\centering
\renewcommand{\arraystretch}{1.3}
\begin{tabular}{p{4cm} p{10cm}}
\toprule
\rowcolor{gray!20}
구분 & 주요 내용 \\
\midrule
기후위기 대응 & 가뭄, 홍수, 폭염 대응, 농업용수 안정성, 재해예방형 기반정비 \\
탄소중립 & 농업부문 온실가스 감축, 재생에너지 확대, 저탄소 농업기술 \\
자원순환ㆍ생태보전 & 물 재이용, 토양ㆍ수질 관리, 생태계 기반 농촌관리 \\
\bottomrule
\end{tabular}
\caption{녹색전환 분석 범위}
\end{table}

\subsection*{해야 할 일}

\begin{itemize}
    \item 녹색전환 관련 핵심 개념 정리
    \begin{itemize}
        \item 기후위기 대응 관련 개념 정리
        \begin{itemize}
            \item 기후변화 적응 (Adaptation)
            \item 회복탄력성 (Resilience)
            \item 취약성 (Vulnerability)
            \item 재해위험관리 (Disaster / Risk Management)
            \item 기후스마트농업 (Climate-Smart Agriculture)
            \item 농업용수 관리 (Agricultural Water Management)
        \end{itemize}
        \item 탄소중립 관련 개념 정리
        \begin{itemize}
            \item 기후변화 완화 (Mitigation)
            \item 탄소중립 (Carbon Neutrality)
            \item 넷제로 (Net Zero)
            \item 탈탄소화 (Decarbonization)
            \item 탄소흡수원 (Carbon Sink)
            \item 재생에너지 (Renewable Energy)
            \item 저탄소농업 (Low-Carbon Agriculture)
        \end{itemize}
        \item 자원순환·생태보전 관련 개념 정리
        \begin{itemize}
            \item 순환경제 (Circular Economy)
            \item 자원순환 (Resource Circulation)
            \item 생태보전 (Ecological Conservation)
            \item 자연기반해법 (Nature-based Solutions)
            \item 생물다양성 (Biodiversity)
            \item 토양·수질 관리 (Soil and Water Management)
            \item 물 재이용 (Water Reuse / Recycling)
        \end{itemize}
    \end{itemize}

    \item 국내외 정책자료 검색을 위한 키워드 리스트 작성
    \begin{itemize}
        \item 국내 정책자료 조사
        \begin{itemize}
            \item 탄소중립 녹색성장 기본계획
            \item 농업·농촌 탄소중립 추진전략
            \item 기후위기 적응대책
            \item 스마트농업 확산 정책
            \item 농업용수 및 기반정비 관련 정책
        \end{itemize}
        \item 해외 정책자료 조사
        \begin{itemize}
            \item EU Green Deal
            \item Farm to Fork Strategy
            \item Common Agricultural Policy (CAP)
            \item USDA Climate-Smart Agriculture
            \item Japan Green Food System Strategy
        \end{itemize}
        \item 국제기구 보고서 및 데이터 검색
        \begin{itemize}
            \item FAO
            \item OECD
            \item World Bank
            \item UNFCCC
            \item IPCC
            \item IEA
        \end{itemize}
        \item 검색 키워드 정리
        \begin{itemize}
            \item climate adaptation agriculture
            \item carbon neutral agriculture
            \item agricultural water management
            \item rural resilience
            \item circular economy agriculture
        \end{itemize}
    \end{itemize}

    \item 분석틀 설정
    \begin{itemize}
        \item 외부 환경 변화 분석
        \begin{itemize}
            \item 기후위기 심화
            \item 글로벌 정책 변화
            \item 기술 발전 및 시장 변화
        \end{itemize}
        \item 정책 변화 분석
        \begin{itemize}
            \item 국내 정책 변화
            \item 해외 정책 변화
        \end{itemize}
        \item 농업·농촌 영향 분석
        \begin{itemize}
            \item 생산 구조 변화
            \item 농촌 생활환경 변화
            \item 기반시설 수요 변화
        \end{itemize}
        \item 농어촌공사 역할 변화 분석
        \begin{itemize}
            \item 기존 사업 재편
            \item 신규 사업 발굴
            \item 중장기 전략 수립
        \end{itemize}
    \end{itemize}

    \item 세부 분석 항목 확정
    \begin{itemize}
        \item 농업 생산 측면
        \begin{itemize}
            \item 생산성 변화
            \item 작목 전환 가능성
            \item 스마트농업 도입
        \end{itemize}
        \item 농촌 인프라 측면
        \begin{itemize}
            \item 농업용수 공급
            \item 배수 및 재해예방 시설
            \item 에너지 인프라
        \end{itemize}
        \item 환경·생태 측면
        \begin{itemize}
            \item 탄소저감 효과
            \item 생태복원
            \item 수질·토양 관리
        \end{itemize}
        \item 정책·제도 측면
        \begin{itemize}
            \item 법·제도 개선
            \item 예산 및 지원체계
            \item ESG 및 녹색금융 연계
        \end{itemize}
        \item 기술·디지털 전환 측면
        \begin{itemize}
            \item 데이터 기반 관리
            \item AI / IoT 활용
            \item 디지털 트윈 / 스마트 인프라
        \end{itemize}
    \end{itemize}
\end{itemize}

\section*{2단계: 대내외 동향 조사}

\subsection*{목적}

2단계에서는 녹색전환과 관련된 국제기구, 주요국, 국내 정부의 정책 동향을 검토하고,
농업ㆍ농촌 부문에 주는 정책적 시사점을 도출한다.

국제적으로는 IPCC, FAO, EU, OECD를 중심으로
기후변화 대응과 식량안보 확보를 위한 정책 전환이 빠르게 진행되고 있다.

국내에서는 2050 탄소중립 정책과 연계하여
농업부문 온실가스 감축, 스마트농업 확대,
통합 물관리 및 재해대응형 기반시설 확충 정책이 추진되고 있다.

\subsection*{해외 동향 조사 대상}

\begin{table}[h]
\centering
\renewcommand{\arraystretch}{1.3}
\begin{tabular}{p{4cm} p{5cm} p{5cm}}
\toprule
\rowcolor{gray!20}
기관ㆍ국가 & 주요 정책ㆍ자료 & 검토 포인트 \\
\midrule
IPCC & 기후변화 평가보고서 & 농업 부문 기후위험, 적응 필요성 \\
FAO & 식량안보 및 농업적응 자료 & 식량안보, 농업 생산성, 기후적응 \\
EU & Green Deal, Farm to Fork & 저탄소 농업, 지속가능 농업체계 \\
OECD & 농촌 회복력 정책 & 농촌지역 회복력, 지역 기반 대응 \\
\bottomrule
\end{tabular}
\end{table}

\section*{3단계: 농업ㆍ농촌 영향 분석}

\subsection*{목적}

3단계에서는 녹색전환이 농업 생산, 자원관리, 재해대응, 식량안보에 미치는 영향을
분야별로 분석한다.

첫째, 생산 측면에서는 기온 상승과 강수 패턴 변화로 인해
작물 생산성 변동성이 증가하고,
재배 가능 지역의 변화가 발생할 수 있다.

둘째, 자원관리 측면에서는 농업용수 부족과 집중호우에 따른
배수 문제의 동시 발생 가능성이 확대된다.

셋째, 재해 대응 측면에서는 가뭄, 홍수, 염해와 같은
기후재해의 빈도와 강도가 증가한다.

넷째, 식량안보 측면에서는 글로벌 공급망 충격과
국내 곡물자급률 한계가 주요 리스크로 작용한다.

\section*{4단계: 국내외 사례 조사}

\subsection*{목적}

4단계에서는 녹색전환에 대응한 국내외 정책 및 사업 사례를 검토하여,
한국농어촌공사에 적용 가능한 시사점을 도출한다.

\begin{table}[h]
\centering
\renewcommand{\arraystretch}{1.3}
\begin{tabular}{p{4cm} p{5cm} p{5cm}}
\toprule
\rowcolor{gray!20}
구분 & 사례 & 검토 포인트 \\
\midrule
해외 & USDA 기후스마트 농업 & 기후적응형 농업지원, 탄소저감 사업 \\
해외 & 네덜란드 스마트 물관리 & 정밀 물관리, 수자원 효율화 \\
해외 & 일본 재해대응 농업기반 & 재해예방형 농업기반시설 관리 \\
국내 & 한국농어촌공사 스마트 물관리 & ICT 기반 농업용수 관리 \\
\bottomrule
\end{tabular}
\end{table}

\section*{5단계: 공사 기능과의 연결}

앞선 분석 결과를 토대로
한국농어촌공사의 기존 기능을 녹색전환 관점에서 재해석한다.

\begin{table}[h]
\centering
\renewcommand{\arraystretch}{1.3}
\begin{tabular}{p{5cm} p{9cm}}
\toprule
\rowcolor{gray!20}
기존 기능 & 녹색전환형 기능 \\
\midrule
농업용수 공급 & 기후적응형 스마트 물관리 \\
농지정비 & 탄소저감형 농지관리 \\
배수개선 & 재해회복형 기반시설 관리 \\
농촌공간개발 & 친환경ㆍ순환형 농촌공간 조성 \\
\bottomrule
\end{tabular}
\end{table}

\section*{6단계: 실천과제 도출}

실천과제는 정책 실행 가능성과 단계별 추진전략을 고려하여
단기, 중기, 장기로 구분하여 도출한다.

\begin{table}[h]
\centering
\renewcommand{\arraystretch}{1.3}
\begin{tabular}{p{3cm} p{11cm}}
\toprule
\rowcolor{gray!20}
구분 & 실천과제 \\
\midrule
단기 & 기후취약지역 진단지도 구축, 재해취약 기반시설 우선 보강 \\
중기 & 탄소중립형 농업기반사업 신설, 농업용수 재이용 시스템 확대 \\
장기 & 기후적응형 농업 인프라 플랫폼 구축 \\
\bottomrule
\end{tabular}
\end{table}

\section*{최종 결과물}

\begin{enumerate}
    \item 녹색전환 동향 분석 문서
    \item 국내외 정책ㆍ사례 비교표
    \item 공사 기능 재정립 표
    \item 단기ㆍ중기ㆍ장기 실천과제 리스트
\end{enumerate}

\section*{정리}

\begin{center}
\textbf{기후위기ㆍ탄소중립ㆍ자원순환 변화}
\\[0.5em]
$\Downarrow$
\\[0.5em]
\textbf{농업ㆍ농촌 생산 및 자원관리 방식 변화}
\\[0.5em]
$\Downarrow$
\\[0.5em]
\textbf{한국농어촌공사의 기능 재정립 및 신규과제 도출}
\end{center}

\end{document}